\clearpage
\section{LSTMs and the limits of recurrent memory}\label{sec:lstm}
\subsection*{Separate stored content from exposed state}
A vanilla RNN replaces its state through a nonlinear transformation at
every step. LSTM introduces an additional cell state and learned gates
that control retaining, writing, and exposing information. The cell state
$c_t$ and hidden state $h_t$ each have width $H$ in the formulation below.

Let $q_t=[e_t;h_{t-1}]\in\R^{d+H}$. The modern cell used in these notes is
\begin{align*}
 f_t&=\sigma(W_fq_t+b_f), & i_t&=\sigma(W_iq_t+b_i),\\
 o_t&=\sigma(W_oq_t+b_o), & g_t&=\tanh(W_gq_t+b_g),\\
 c_t&=f_t\odot c_{t-1}+i_t\odot g_t,
 &h_t&=o_t\odot\tanh(c_t).
\end{align*}
Here the four gate/candidate matrices have shape $H\times(d+H)$.
Within this cell, $W_o,b_o$ name the output gate; a separate vocabulary
readout acts on $h_t$ to predict the next token.

\fig{lstm-cell}{A cell with a forget gate and no peephole connections. The top
path carries the cell state; multiplication nodes apply gates elementwise.}

\begin{itemize}
\item The \textbf{forget gate} $f_t$ controls how much previous cell content survives.
\item The \textbf{input gate} $i_t$ controls how much candidate content $g_t$ is written.
\item The \textbf{output gate} $o_t$ controls how much transformed cell content is exposed as $h_t$.
\end{itemize}
Sigmoid gates lie between zero and one, independently for each coordinate.
They are not a categorical distribution: $f_t+i_t$ need not equal one.
The candidate can be negative. The cell state accumulates content through
an additive update and is not itself restricted to $[-1,1]$.

\textbf{Historical convention.} \readingref{R10}, the 1997 LSTM paper,
introduced the memory-cell idea before this forget-gate formulation.
The extension is documented by \href{https://sferics.idsia.ch/pub/juergen/FgGates-NC.pdf}{Gers,
Schmidhuber, and Cummins (2000)}. Distinguishing the two avoids attributing
all modern equations to the original cell.

\selfcheck{Why keep $c_t$ as well as $h_t$? The model can retain cell content
while controlling how much of it reaches the prediction and next gate computations.}

\nextpage{A scalar LSTM update and its gradient path}
Choose $c_{t-1}=0.8$, $f_t=0.9$, $i_t=0.2$, $g_t=0.5$, and $o_t=0.75$.
These hand-chosen values make each operation visible:
\[
 c_t=0.9(0.8)+0.2(0.5)=0.82,\qquad
 h_t=0.75\tanh(0.82)\approx0.5063.
\]
The previous memory contributes $0.72$ and the candidate contributes $0.10$.
The hidden state exposes a bounded transform of their sum.

\fig{gate-arithmetic}{Idealized gate settings illustrate retain, overwrite,
and hide operations. Exact zero/one gates are limiting cases of sigmoid gates.}

If the gates and candidate are held fixed, increasing $c_{t-1}$ by $0.1$
increases $c_t$ by $0.09$. Along this direct cell-state path,
\[
 \left.\frac{\partial c_t}{\partial c_{t-1}}\right|_{f_t,i_t,g_t\ \mathrm{fixed}}
 =\operatorname{diag}(f_t).
\]
Across several steps, the contribution through only the direct cell path
is an elementwise product of forget gates. Gates near one can preserve
this contribution better than repeated saturated nonlinear transformations.

\subsection*{Why this is not the full gradient}
The gates depend on $h_{t-1}$, which depends on earlier cells and gates.
There are therefore additional paths from earlier states to later losses.
A product of forget gates isolates one path; it does not equal the full
recurrent derivative in general. The output gate and $\tanh(c_t)$ also
affect the sensitivity of a prediction to the stored content.

Even on the direct path, a forget value of 0.9 repeated 50 times yields
about 0.00515. Gates closer to one help retention, but successful learning
also requires that the right content be written, preserved, and exposed
at the right time. An LSTM improves the architecture's options for learning
memory; it does not eliminate all gradient or optimization difficulties.

\textbf{A useful alternative.} A GRU combines state and gating differently,
using reset and update gates. \readingref{R12} provides the historical
encoder--decoder context for that design. For this recap, understand the
shared purpose of controlling recurrent memory before comparing full equations.

\takeaway{Treat the additive memory path as a useful mechanism with stated
assumptions. Separate its local derivative from the total gradient of the network.}
\selfcheck{Can $h_t$ be near zero while $c_t$ retains a nonzero value? Yes:
a nearly closed output gate can hide stored content from the exposed state.}

\nextpage{What recurrent memory still cannot make easy}
\subsection*{Sequential computation remains}
The standard RNN or LSTM update needs the previous state. Even when all
training tokens are known, state $h_t$ cannot be evaluated by the usual
recurrence until $h_{t-1}$ is available. Matrix operations within a step
and different sequences in a batch can run in parallel; the within-sequence
chain remains a source of limited parallelism.

\fig{sequential-computation}{Two sequences can be processed alongside each
other, but each row has a left-to-right state dependency. Knowing the next
input does not supply the next hidden state.}

\subsection*{A state must retain the distinctions needed later}
A fixed-width state can be sufficient for some tasks, such as tracking a
small set of conditions. It need not preserve every detail of the prefix.
However, when useful information is spread across a long sequence, later
predictions depend on what survived earlier updates. Long-range credit
assignment must also teach the network which details to retain.

\fig{encoder-bottleneck}{A classic encoder--decoder interface passes only a
final encoder representation to the decoder. This particular design requires
that representation to carry the source information needed for every output.
Decoder token inputs are omitted.}

In the sequence-to-sequence model of \readingref{R13}, an LSTM encoder
produces a fixed-dimensional representation for an LSTM decoder. Unlike
an ordinary next-token LM on one stream, the task maps one sequence to
another. The interface makes the burden concrete: details needed at
different output positions must all remain available through that final
representation. This bottleneck is a property of the interface, not a
claim that every recurrent architecture exposes only one source vector.

The remaining concerns are related but distinct: sequential computation,
learning across long gradient paths, and retaining useful distinctions in
a finite state. Gates address part of the memory problem; they do not
remove the serial dependency or automatically change the encoder--decoder interface.

\takeaway{The next question is how a prediction could consult earlier
information more directly. Lecture 05 begins with learned alignment and
attention, using \readingref{R14} and \readingref{R15}.}
\selfcheck{Which limitation would remain if a recurrent model remembered
every needed detail perfectly? Its standard state updates would still depend on earlier steps.}
